\documentclass[10pt,a4paper]{amsart}
\usepackage{amsmath,amssymb,amsthm,graphicx}
% Retained minimal preamble for the Stacks Project chapter more-morphisms.tex (More on Morphisms).
% Verbatim excerpt of the chapter preamble: Xy-pic setup (native preamble lines 14-18) and the native
% theorem declarations with the mathematical macros (native preamble lines 154-217). Copied mechanically;
% nothing added, removed or redefined. The wrapper class and the amsmath/amssymb/amsthm/graphicx base
% packages are supplied by the assembler, so the class/externaldocument part of the native preamble is omitted.
% For commutative diagrams we use Xy-pic
\usepackage[all]{xy}

% We use 2cell for 2-commutative diagrams.
\xyoption{2cell}
% Theorem environments.
%
\theoremstyle{plain}
\newtheorem{theorem}[subsection]{Theorem}
\newtheorem{proposition}[subsection]{Proposition}
\newtheorem{lemma}[subsection]{Lemma}

\theoremstyle{definition}
\newtheorem{definition}[subsection]{Definition}
\newtheorem{example}[subsection]{Example}
\newtheorem{exercise}[subsection]{Exercise}
\newtheorem{situation}[subsection]{Situation}

\theoremstyle{remark}
\newtheorem{remark}[subsection]{Remark}
\newtheorem{remarks}[subsection]{Remarks}

\numberwithin{equation}{subsection}

% Macros
%
\def\lim{\mathop{\mathrm{lim}}\nolimits}
\def\colim{\mathop{\mathrm{colim}}\nolimits}
\def\Spec{\mathop{\mathrm{Spec}}}
\def\Hom{\mathop{\mathrm{Hom}}\nolimits}
\def\Ext{\mathop{\mathrm{Ext}}\nolimits}
\def\SheafHom{\mathop{\mathcal{H}\!\mathit{om}}\nolimits}
\def\SheafExt{\mathop{\mathcal{E}\!\mathit{xt}}\nolimits}
\def\Sch{\mathit{Sch}}
\def\Mor{\mathop{\mathrm{Mor}}\nolimits}
\def\Ob{\mathop{\mathrm{Ob}}\nolimits}
\def\Sh{\mathop{\mathit{Sh}}\nolimits}
\def\NL{\mathop{N\!L}\nolimits}
\def\CH{\mathop{\mathrm{CH}}\nolimits}
\def\proetale{{pro\text{-}\acute{e}tale}}
\def\etale{{\acute{e}tale}}
\def\QCoh{\mathit{QCoh}}
\def\Ker{\mathop{\mathrm{Ker}}}
\def\Im{\mathop{\mathrm{Im}}}
\def\Coker{\mathop{\mathrm{Coker}}}
\def\Coim{\mathop{\mathrm{Coim}}}

% Boxtimes
%
\DeclareMathSymbol{\boxtimes}{\mathbin}{AMSa}{"02}

%
% Macros for moduli stacks/spaces
%
\def\QCohstack{\mathcal{QC}\!\mathit{oh}}
\def\Cohstack{\mathcal{C}\!\mathit{oh}}
\def\Spacesstack{\mathcal{S}\!\mathit{paces}}
\def\Quotfunctor{\mathrm{Quot}}
\def\Hilbfunctor{\mathrm{Hilb}}
\def\Curvesstack{\mathcal{C}\!\mathit{urves}}
\def\Polarizedstack{\mathcal{P}\!\mathit{olarized}}
\def\Complexesstack{\mathcal{C}\!\mathit{omplexes}}
% \Pic is the operator that assigns to X its picard group, usage \Pic(X)
% \Picardstack_{X/B} denotes the Picard stack of X over B
% \Picardfunctor_{X/B} denotes the Picard functor of X over B
\def\Pic{\mathop{\mathrm{Pic}}\nolimits}
\def\Picardstack{\mathcal{P}\!\mathit{ic}}
\def\Picardfunctor{\mathrm{Pic}}
\def\Deformationcategory{\mathcal{D}\!\mathit{ef}}
\newenvironment{reference}{}{}
\newenvironment{slogan}{}{}
\newenvironment{history}{}{}
\usepackage{cleveref}
\crefname{originallemma}{original lemma}{original lemmas}
\Crefname{originallemma}{Original lemma}{Original lemmas}
\makeatletter
\newlabel{coherent-lemma-ample-in-neighbourhood}{{0D2N}{source}{}{sourcepointer-1}{}}
\newlabel{coherent-lemma-ample-in-neighbourhood@cref}{{[originallemma][0][]0D2N}{[1][0][]source}{}{}{}}
\makeatother
\crefname{originallemma}{original lemma}{original lemmas}
\Crefname{originallemma}{Original lemma}{Original lemmas}
\makeatletter
\newlabel{limits-lemma-descend-invertible-modules}{{0B8W}{source}{}{sourcepointer-2}{}}
\newlabel{limits-lemma-descend-invertible-modules@cref}{{[originallemma][0][]0B8W}{[1][0][]source}{}{}{}}
\makeatother
\crefname{originallemma}{original lemma}{original lemmas}
\Crefname{originallemma}{Original lemma}{Original lemmas}
\makeatletter
\newlabel{limits-lemma-limit-ample}{{09MT}{source}{}{sourcepointer-3}{}}
\newlabel{limits-lemma-limit-ample@cref}{{[originallemma][0][]09MT}{[1][0][]source}{}{}{}}
\makeatother
\crefname{originallemma}{original lemma}{original lemmas}
\Crefname{originallemma}{Original lemma}{Original lemmas}
\makeatletter
\newlabel{limits-lemma-proper-limit-of-proper-finite-presentation-noetherian}{{0A0P}{source}{}{sourcepointer-4}{}}
\newlabel{limits-lemma-proper-limit-of-proper-finite-presentation-noetherian@cref}{{[originallemma][0][]0A0P}{[1][0][]source}{}{}{}}
\makeatother
\crefname{originallemma}{original lemma}{original lemmas}
\Crefname{originallemma}{Original lemma}{Original lemmas}
\makeatletter
\newlabel{morphisms-lemma-ample-base-change}{{0893}{source}{}{sourcepointer-5}{}}
\newlabel{morphisms-lemma-ample-base-change@cref}{{[originallemma][0][]0893}{[1][0][]source}{}{}{}}
\makeatother
\crefname{originallemma}{original lemma}{original lemmas}
\Crefname{originallemma}{Original lemma}{Original lemmas}
\makeatletter
\newlabel{morphisms-lemma-pullback-ample-tensor-relatively-ample}{{0892}{source}{}{sourcepointer-6}{}}
\newlabel{morphisms-lemma-pullback-ample-tensor-relatively-ample@cref}{{[originallemma][0][]0892}{[1][0][]source}{}{}{}}
\makeatother
\title{Ampleness of a line bundle on one fibre of a proper morphism spreads to an open neighbourhood of the base point}
\author{The Stacks Project Authors}
\date{stacks-project:more-morphisms.tex:local-snapshot-20261011; GNU FDL 1.2 or later}
\begin{document}
\maketitle
\noindent\textit{Editorial note.} Statement span is lines 14153-14163 of sources/stacks-more-morphisms/pinned-original-source.tex; proof span is lines 14165-14194, the complete author proof (28 body lines).\par\medskip
\noindent\textit{Editorial note.} Navigation stubs: 6 cited-but-not-included labels are supplied as non-load-bearing pointers with their pinned tags where the snapshot tag index has them.\par\medskip
\noindent\textit{Editorial note.} extra\_packages is empty on purpose: the retained preamble already loads xy with options and a second declaration clashes in the assembler wrapper.\par\medskip
\noindent\textit{Editorial note.} Text is retained verbatim, including the author's own typos and placeholders.\par\medskip
\noindent\textit{Editorial note.} The proof closes with author-supplied hints rather than a line-by-line verification of the last pullback; the mechanism (affine base change preserving ampleness) is fully named in the span. Reference block cites \textbackslash{}cite[IV Corollary 9.6.4]\{EGA\}; bibliography.tex with the snapshot my.bib EGA entry attached.\par\medskip
\noindent\textit{Editorial note.} The statement's reference environment cites \textbackslash{}cite[IV Corollary 9.6.4]\{EGA\}; a bibliography.tex holding the snapshot my.bib EGA entry is attached (precedent L16058/L16018).\par\medskip
\noindent\textit{Allowed standard prerequisites (editorial).}\begin{itemize}
\item Explicit prerequisite boundary. The Noetherian approximation of proper maps (Limits, proposition and lemma-proper-limit-*), descent of invertible modules to a finite stage (Limits, lemma-descend-invertible-modules), the limit-detection of ampleness (Limits, lemma-limit-ample), and the finite-type openness of ampleness (Cohomology of Schemes, lemma-ample-in-neighbourhood).
\end{itemize}\medskip
\ifdefined\hypertarget\hypertarget{sourcepointer-1}{}\fi
\noindent\textit{Source reference.} Original lemma 0D2N: sources/stacks-more-morphisms/pinned-original-source.tex, source lines 14194--14194. Cited by the retained statement or proof but not part of the collected spans; supplied only as a navigation pointer. The argument does not depend on its content.\par\medskip
\ifdefined\hypertarget\hypertarget{sourcepointer-2}{}\fi
\noindent\textit{Source reference.} Original lemma 0B8W: sources/stacks-more-morphisms/pinned-original-source.tex, source lines 14194--14194. Cited by the retained statement or proof but not part of the collected spans; supplied only as a navigation pointer. The argument does not depend on its content.\par\medskip
\ifdefined\hypertarget\hypertarget{sourcepointer-3}{}\fi
\noindent\textit{Source reference.} Original lemma 09MT: sources/stacks-more-morphisms/pinned-original-source.tex, source lines 14194--14194. Cited by the retained statement or proof but not part of the collected spans; supplied only as a navigation pointer. The argument does not depend on its content.\par\medskip
\ifdefined\hypertarget\hypertarget{sourcepointer-4}{}\fi
\noindent\textit{Source reference.} Original lemma 0A0P: sources/stacks-more-morphisms/pinned-original-source.tex, source lines 14194--14194. Cited by the retained statement or proof but not part of the collected spans; supplied only as a navigation pointer. The argument does not depend on its content.\par\medskip
\ifdefined\hypertarget\hypertarget{sourcepointer-5}{}\fi
\noindent\textit{Source reference.} Original lemma 0893: sources/stacks-more-morphisms/pinned-original-source.tex, source lines 14194--14194. Cited by the retained statement or proof but not part of the collected spans; supplied only as a navigation pointer. The argument does not depend on its content.\par\medskip
\ifdefined\hypertarget\hypertarget{sourcepointer-6}{}\fi
\noindent\textit{Source reference.} Original lemma 0892: sources/stacks-more-morphisms/pinned-original-source.tex, source lines 14194--14194. Cited by the retained statement or proof but not part of the collected spans; supplied only as a navigation pointer. The argument does not depend on its content.\par\medskip
\begin{lemma}
\label{lemma-ample-in-neighbourhood}
\begin{reference}
\cite[IV Corollary 9.6.4]{EGA}
\end{reference}
Let $f : X \to Y$ be a proper morphism of schemes.
Let $\mathcal{L}$ be an invertible $\mathcal{O}_X$-module.
Let $y \in Y$ be a point such that $\mathcal{L}_y$ is ample
on $X_y$. Then there is an open neighbourhood $V \subset Y$
of $y$ such that $\mathcal{L}|_{f^{-1}(V)}$ is ample on $f^{-1}(V)/V$.
\end{lemma}

\begin{proof}
We may assume $Y$ is affine. Then we find a directed set $I$
and an inverse system of morphisms $X_i \to Y_i$ of schemes
with $Y_i$ of finite type over $\mathbf{Z}$, with affine
transition morphisms $X_i \to X_{i'}$ and $Y_i \to Y_{i'}$,
with $X_i \to Y_i$ proper, such that $X \to Y = \lim (X_i \to Y_i)$.
See Limits, Lemma
\ref{limits-lemma-proper-limit-of-proper-finite-presentation-noetherian}.
After shrinking $I$ we can assume we have a compatible system of
invertible $\mathcal{O}_{X_i}$-modules $\mathcal{L}_i$
pulling back to $\mathcal{L}$, see
Limits, Lemma \ref{limits-lemma-descend-invertible-modules}.
Let $y_i \in Y_i$ be the image of $y$.
Then $\kappa(y) = \colim \kappa(y_i)$.
Hence for some $i$ we have $\mathcal{L}_{i, y_i}$
is ample on $X_{i, y_i}$ by
Limits, Lemma \ref{limits-lemma-limit-ample}.
By Cohomology of Schemes, Lemma \ref{coherent-lemma-ample-in-neighbourhood}
we find an open neighbourhood
$V_i \subset Y_i$ of $y_i$ such that
$\mathcal{L}_i$ restricted to $f_i^{-1}(V_i)$
is ample relative to $V_i$.
Letting $V \subset Y$ be the inverse image of
$V_i$ finishes the proof (hints: use
Morphisms, Lemma \ref{morphisms-lemma-ample-base-change} and
the fact that $X \to Y \times_{Y_i} X_i$ is affine
and the fact that the pullback of an
ample invertible sheaf by an affine morphism is ample by
Morphisms, Lemma \ref{morphisms-lemma-pullback-ample-tensor-relatively-ample}).
\end{proof}

\begin{thebibliography}{1}
\bibitem{EGA}
Dieudonn{\'e}, Jean and Grothendieck, Alexander.
\newblock {\em {\'E}l{\'e}ments de g{\'e}om{\'e}trie alg{\'e}brique}.
\newblock Inst. Hautes {\'E}tudes Sci. Publ. Math. 4, 8, 11, 17, 20, 24, 28, 32 (1961--1967).
\end{thebibliography}
\end{document}
